\documentclass[11pt,a4paper]{article}
\usepackage[margin=27mm]{geometry}
\usepackage{amsmath,amssymb,amsthm,booktabs,hyperref}
\newtheorem{lemma}{Lemma}
\newtheorem{theorem}[lemma]{Theorem}
\newtheorem{corollary}[lemma]{Corollary}
\newcommand{\Sz}{\operatorname{Sz}}
\newcommand{\C}{\operatorname{C}}
\newcommand{\F}{\mathbb F}
\newcommand{\gen}[1]{\langle#1\rangle}
\title{The necessary patterns when the split torus has prime order}
\author{Human Compression project}\date{}
\begin{document}\maketitle
\begin{abstract}
We prove the necessary-list assertion of Theorem 5.3 in the fixed 1999
rank-three Suzuki paper. Proper covers of a subgroup of prime order
reduce the possibilities to three patterns. A sparse matrix entry excludes
one apparent additional configuration without assuming an unstated
intersection property of ovoid circles. The prime-degree rank bound and
the rank-two locally doubly transitive cases follow as well. We distinguish
the source's actions on residues defined by subgroup intersections from
actions on full incidence residues. Under the latter hypothesis we also
prove geometric thinness; it does not follow merely by identifying the
former actions.
\end{abstract}

\section{Definitions and acquired restrictions}
Let $q=2^n\ge8$, suppose $p=q-1$ is prime, and put $G=\Sz(q)$.
Then $n$ is prime. We use the source's subgroup condition: for every
proper $J\subsetneq I$, some $G_{J\cup\{i\}}$ is maximal proper in
$G_J=\bigcap_{j\in J}G_j$, with $G_\varnothing=G$.
The preceding note \emph{A rank obstruction in prime-degree Suzuki groups}
proves directly from this condition that every additional type strictly
decreases an intersection. It excludes a point stabilizer as a parabolic
in every rank at least two, and a full nonsplit torus normalizer in every
rank at least three. We import those arguments; they do not assume flag
transitivity. Their distinct mathematical review has been adopted in the accompanying
prime-degree root review; the exact accepted inputs are bound below.

We retain the same classical Suzuki subgroup classification, torus
partition, ovoid action and faithful matrix model. A split torus $A=C_p$
fixes exactly two ovoid points, denoted $P,Q$; its normalizer is
$D=A\rtimes\gen w$, where $w$ interchanges $P,Q$ and inverts $A$.
Write $Z_P,Z_Q$ for the centres of their Sylow two-subgroups. Every
nonidentity involution has one fixed ovoid point, and its centralizer is
the corresponding Sylow two-subgroup. There are no proper simple subfield
groups because $n$ is prime.

\begin{lemma}\label{covers}
The following proper overgroups in $G$ have the indicated forms.
\begin{enumerate}
\item If $C_2$ is maximal in $L<G$, then $L$ has order four or
      $L=D_{2r}$ for an odd prime $r$ dividing $|G|$.
\item The proper overgroups of a full split torus $A=C_p$ are
      $A,D,Z_PA,Z_QA,B_P,B_Q$, where $B_P,B_Q$ are the two point
      stabilizers containing $A$.
\item If $C_r$ is maximal in $L<G$ for an odd prime $r$, and $L$
      contains an involution, then $L=D_{2r}$ or $r=p$ and
      $L=Z_RA'$ for a split torus $A'$ of order $p$ fixing $R$.
\end{enumerate}
\end{lemma}
\begin{proof}
For (1), first suppose $L$ lies in a point stabilizer $B=U\rtimes T$.
Its normal two-subgroup $V=L\cap U$ contains the given $C_2$.
Maximality gives $V=L$ or $V=C_2$. In the latter case $L$ normalizes
$C_2$ and hence centralizes its involution, so again $L$ is a two-group.
A maximal subgroup of a finite two-group has index two, giving $|L|=4$.
If $L$ lies in a cyclic torus normalizer, intersect it with the odd
cyclic kernel. A nontrivial intersection $K$ is normal and
$K C_2>C_2$, so maximality forces $L=K C_2$. The involution inverts
$K$. Its order must be prime, since otherwise a proper nontrivial
rotation subgroup gives an intermediate dihedral subgroup. If $K=1$,
$L$ is a two-group and the same index argument applies. These cases
exhaust the proper subgroups by the classical classification.

For (2), coprimality excludes the nonsplit normalizers. A split
normalizer containing $A$ is $D$, by uniqueness of the torus. A point
stabilizer containing $A$ fixes $P$ or $Q$. In either one, write an
overgroup $L$ of $A$ as $VA$, with $V=L\cap U$. The action of $A$
on both $Z(U)$ and $U/Z(U)$ is the full scalar action of $\F_q^*$,
so is irreducible over $\F_2$. If $V\le Z(U)$, then $V$ is zero
or the whole centre. Otherwise $VZ(U)=U$, and the acquired identity
$Z(U)=\Phi(U)$ forces $V=U$. This gives precisely the displayed list.

For (3), a point-stabilizer case has $r\mid p$, so $r=p$ and
(2) shows that the only cover there is $Z_RA'$. In a cyclic torus
normalizer, let $K$ be the intersection of $L$ with the odd kernel.
It contains $C_r$ and is proper in $L$, because $L$ has an involution.
Maximality gives $K=C_r$. The quotient embeds in a cyclic group of
order two or four. An order-four quotient would give an intermediate
subgroup, so the quotient has order two. Its involution acts by
inversion on the odd torus. Thus $L=D_{2r}$.
\end{proof}

\section{An excluded triangle}
\begin{lemma}\label{triangle}
There is no family satisfying the source's subgroup condition with
\[
 G_0=D,\qquad G_{01}=A,\qquad G_{02}\cong C_2,
 \qquad G_{12}\cong C_p,\qquad G_{012}=1.
\]
\end{lemma}
\begin{proof}
By Lemma~\ref{covers}(2) and the point-stabilizer exclusion,
$G_1$ is $Z_PA$ or $Z_QA$. Interchanging $P,Q$ if necessary makes
$G_1=Z_PA$. Conjugation by $A$, which preserves $G_1$, sends the
reflection in $G_{02}$ to $w$: squaring is bijective on the odd group
$A$. Every order-$p$ subgroup of $Z_PA$ is a complement and is conjugate
to $A$ by $Z_P$. This is the elementary affine fixed-point argument
already used in the accepted preliminary note. Thus
$G_{12}=A^z$ for some $1\ne z\in Z_P$; $z=1$ would contradict
$G_{012}=1$.

Apply Lemma~\ref{covers}(2) to $A^z$. The group $G_2$ is a proper
overgroup containing $w$. Point stabilizers are excluded and $G_2$ cannot
equal $A^z$, so the possibilities are $D^z$ or one of the two groups
$Z_RA^z$ at its fixed points. Those fixed points are $P,Q^z$.
The group at $P$ is $G_1$, which is excluded by strict intersections.

The possibility $w\in D^z$ is also impossible. Such a $w$ would
interchange $P,Q^z$, whereas $w$ interchanges $P,Q$. Hence $Q^z=Q$.
The nonidentity involution $z$ fixes only $P$, a contradiction.
For the remaining possibility, $w\in B_{Q^z}$ would give
$zwz\in B_Q$, since $z^2=1$. Use the retained faithful matrix model
with $B_P$ lower triangular, $B_Q$ upper triangular, $w$ antidiagonal,
and
\[
 z=z(b)=
 \begin{pmatrix}
 1&0&0&0\\0&1&0&0\\b&0&1&0\\b^\theta&b&0&1
 \end{pmatrix},\qquad
 b\in\F_q^*,\quad\theta=2^{(n+1)/2}.
\]
Direct multiplication gives $(z(b)wz(b))_{21}=b\ne0$. It cannot
be upper triangular. This excludes the last possibility.
\end{proof}

The matrix calculation is a quantified identity in the field; it is
not a finite test. It replaces the source's ovoid-orbit intersection
step in this particular configuration. The matrix model itself is
classical and is the same one used in the retained Suzuki realization
manuscripts.

\section{All necessary rank-three patterns}
\begin{theorem}\label{patterns}
Every rank-three subgroup family satisfying the source's condition has,
after permuting its types, one of the following intersection patterns:
\[
\begin{array}{c|c|c|c|c|c}
 G_0&G_1&G_2&G_{01}&G_{02}&G_{12}\\\hline
 D_{2p}&E_q\rtimes C_p&E_q\rtimes C_p&C_2&C_2&C_p\\
 D_{2p}&E_q\rtimes C_p&D_{2r}&C_p&C_2&C_2\\
 D_{2p}&D_{2r}&D_{2s}&C_2&C_2&C_2
\end{array}
\]
The triple intersection is trivial. In the second row $r$ is an odd
prime divisor of $(q-1)(q^2+1)$. In the third row $s$ is such an odd
prime, and $r$ is either such an odd prime or $2$; the notation
$D_4$ means $C_2^2$. Equal odd primes are allowed.
These are precisely the necessary patterns of source Theorem 5.3.
\end{theorem}
\begin{proof}
Some parabolic is maximal in $G$. The imported restrictions exclude
point stabilizers and full nonsplit normalizers, leaving $G_0=D$.
Inside $D_{2p}$ the proper nontrivial subgroups are $A=C_p$ and its
reflection groups $C_2$. Strict intersections and the subgroup condition
therefore leave exactly two possibilities for its rank-two residue:
$(G_{01},G_{02},G_{012})=(A,C_2,1)$ or $(C_2,C_2,1)$.
The triple intersection is maximal in $G_{12}$ as well, so $G_{12}$
is cyclic of prime order.

In the first case, Lemma~\ref{covers}(2), strict intersections and the
point-stabilizer exclusion give $G_1=Z_PA$, after exchanging the two
fixed points. The prime order of $G_{12}$ is either $2$ or $p$.
Lemma~\ref{triangle} excludes $p$. Thus $G_2$ contains two distinct
involutions, with one of their order-two groups maximal in $G_2$ by
the subgroup condition at $\{2\}$. Lemma~\ref{covers}(1) gives an
odd prime dihedral group or a four-group; a cyclic group of order four
has only one involution. The four-group possibility is impossible:
the involution in $G_{12}\le Z_P$ fixes $P$, whereas the reflection
in $G_{02}\le D$ exchanges $P,Q$. Commuting involutions would lie
in the centralizer Sylow two-subgroup of the first and would both
fix $P$. We obtain the second row.

In the second residue case, write $G_{12}=C_r$. If $r=2$, each of
$G_1,G_2$ contains two distinct involutions and has an order-two
subgroup maximal in it. Lemma~\ref{covers}(1) makes each a prime
dihedral group or a four-group. They cannot both be four-groups:
if $u,v$ are the two distinct reflections in $G_0$ and $z$ is the
involution of $G_{12}$, both $u,v$ would centralize $z$. Then
$\gen{u,v}=D_{2p}$ would lie in the two-group $\C_G(z)$, impossible.
This gives exactly the third row.

Suppose instead that $r$ is odd. At each $G_i$, $i=1,2$, either
its intersection $C_2$ with $G_0$ or its intersection $C_r$ with the
other parabolic is maximal. Lemma~\ref{covers}(1),(3) gives
$G_i=D_{2r}$, or $r=p$ and $G_i$ is $E_q\rtimes C_p$.
If $r\ne p$, both are $D_{2r}$ and lie in the unique nonsplit
normalizer determined by $G_{12}$. That normalizer would contain
the two reflections generating $D_{2p}$, contradicting coprimality
of its odd order with $p$. Thus $r=p$.
If both are $D_{2p}$, their common rotation subgroup $G_{12}$
determines the same full split normalizer, so they coincide, contrary
to strict intersections. If exactly one is dihedral, permuting types
gives the excluded triangle of Lemma~\ref{triangle}. Both must therefore
be $E_q\rtimes C_p$, giving the first row.
\end{proof}

This is a necessity proof. The previously reviewed phase-five
realizations show sufficiency at the level of the specified subgroup
intersections, but are not used anywhere in this proof. Their additional
geometric scope must be stated separately when imported.

\begin{corollary}\label{rank}
Under the same prime-$q-1$ hypothesis no such family has rank at least four.
\end{corollary}
\begin{proof}
A maximal parabolic must again be $D_{2p}$. Two successive applications
of the subgroup condition give a maximal proper subgroup of $D_{2p}$,
then a maximal proper subgroup of that group. The first has prime order
$p$ or $2$, and the second is trivial. If rank is at least four,
the condition at this three-type intersection would demand a proper
maximal subgroup of the trivial group. This is impossible.
\end{proof}

\section{Local actions and incidence panels}
On printed page~45, the source defines $\Gamma_F$ for a flag
$F\subseteq C=\{G_0,G_1,G_2\}$ by the intersection subgroups.
On printed page~46 its local double transitivity condition refers to
flags $F\subset C$ and to this $\Gamma_F$, and it expressly qualifies
agreement with the earlier definition in its reference~[10] by flag
transitivity. We preserve that literal convention and keep the full
incidence panel separate. For a two-type base flag of types $i,j$,
the algebraic rank-one residue is $G_{ij}/G_{012}$.
In the three patterns above its denominator is trivial, so its action
is regular. A regular action of a nontrivial group is doubly transitive
precisely when the group has order two: its point stabilizer is trivial,
hence is transitive on the remaining points precisely when there is only
one such point. Thus the third row, and only that row, has all the
source's algebraic rank-one actions doubly transitive.

The corresponding geometric implication under the stronger interpretation
can be proved without any assumption of flag transitivity.
\begin{lemma}\label{actual}
Suppose a rank-three family in Theorem~\ref{patterns} has the property
that the stabilizer of each two-type base flag is doubly transitive on
its full incidence panel. Then the geometry is thin and flag transitive.
\end{lemma}
\begin{proof}
For the base flag $\{G_i,G_j\}$, its stabilizer is $G_{ij}$.
The orbit of the remaining base coset $G_k$ has stabilizer
$G_{012}=1$ and therefore has size $|G_{ij}|$. Transitivity on the
full panel forces this orbit to be the entire panel. Consequently the
action on it is regular, and double transitivity forces $|G_{ij}|=2$.
Every two-type flag is a translate of the base flag of its type:
the two incident cosets share a representative. All panels therefore
have exactly two elements, proving thinness. Moreover every chamber
can first be translated so that two of its elements are base cosets;
transitivity of their stabilizer on the panel then moves its third
element to the remaining base coset. This proves flag transitivity.
\end{proof}

For rank two, the algebraic residue agrees with the actual incidence
residue. The accepted complete normalizer families give a short
independent check of the remaining prime branches of Theorem~6.1.
The split-normalizer family with intersection $C_p$ has local actions
of degrees two and $q$; the second is the natural affine action of
$E_q\rtimes\F_q^*$ and is doubly transitive. This is one equivalence
class. Every other split-normalizer family has intersection $C_2$,
and its action on $p\ge7$ cosets is not doubly transitive: a point
stabilizer of order two cannot be transitive on $p-1>2$ points.

In a nonsplit family the maximal parabolic is $C_t\rtimes C_4$,
where $t=a_\epsilon(q)$ is prime, and the intersection is $C_4$.
Its action on $t$ cosets is doubly transitive exactly when $t=5$.
Indeed the point stabilizer has order four, acts without fixed
nonzero vectors in the cyclic kernel, and is transitive there exactly
when $t-1=4$. The equality $a_\epsilon(q)=5$ occurs only for $q=8$
and the minus sign: it holds there, whereas $q\ge32$ gives
$a_\epsilon(q)\ge q/2+1\ge17$.
At $q=8$ the other parabolic must be another $C_5\rtimes C_4$ or
$C_4\times C_2$. Both local actions are doubly transitive (degree
five or degree two). Their accepted class counts are each
$(q-2)/(2n)=1$. The opposite-torus family has an action of degree
thirteen and fails the same test. Hence there is exactly one rank-two
class when $q>8$, and exactly three when $q=8$, in the source's
specified-group equivalence convention.

\paragraph{The source's thinness conclusion.}
Under the literal algebraic definition, the source further calls all
these rank-three locally doubly transitive cases thin. The two cosets
in an algebraic residue need not a priori exhaust the third-type elements
incident with a pair of cosets in the full incidence structure.
Lemma~\ref{actual} proves thinness if local double transitivity is
required on the full panel; it does not identify the two hypotheses.
The retained diagnostic in the next section checks one $\Sz(8)$ marking
and finds exactly two elements in each panel. It does not distinguish
the two hypotheses in that example and is not a counterexample. The
literal source's general thinness conclusion remains a separate obligation.

\section{One retained marking over the field of eight elements}
The sole finite acquisition in this component revisits one previously
constructed marking. It searches no additional marking and enumerates
no ambient Suzuki group. Put
\[
 \F_8=\F_2[t]/(t^3+t+1),\quad
 a=w,\quad b=z(t)wz(t),\quad c=z(1),
\]
\[
 G_0=\gen{a,b},\qquad G_1=\gen{a,c},\qquad G_2=\gen{b,c},
\]
where $w$ is the antidiagonal matrix and $z(b)$ is the matrix in
Lemma~\ref{triangle}, now with $\theta=4$.
The returned certificate gives the complete element lists of the three
subgroups, of orders $14,10,10$. An independent reader implements
multiplication in this displayed field, verifies every subgroup product
and inverse, and checks that the indicated generators produce the lists.
Their pair intersections have order two and their triple intersection
is trivial. The first subgroup is the full split torus normalizer in
$\Sz(8)$ by the classical subgroup description; each of its two distinct
involutions inverts its cyclic subgroup of order seven. The three
parabolics therefore satisfy the source's subgroup condition: $G_0$ is
maximal in $G$, each rank-one intersection is maximal in either incident
prime dihedral parabolic, and the trivial group is maximal in every
order-two intersection. This statement uses the retained classical
Suzuki model and maximal-subgroup classification, not a new ambient census.

Here is the complete test of actual panels. Fix the base flag
$\{G_j,G_k\}$ and let $i$ be the missing type. A left coset of $G_i$
incident with $G_j$ has a representative in $G_j$, because incidence
means nonempty intersection. Form the distinct sets $gG_i$ for all
$g\in G_j$, and retain exactly those meeting $G_k$. This gives the
entire panel, with no assumption of flag transitivity. Such a coset
belongs to the base chamber's $G_j\cap G_k$ orbit precisely when it
meets $G_j\cap G_k$. The actual returned values are
\[
\begin{array}{c|c|c|c}
 \text{missing type}&\text{candidate cosets}&\text{actual panel}&
 \text{cosets in the standard orbit}\\\hline
 0&5&2&2\\1&7&2&2\\2&7&2&2.
\end{array}
\]
The certificate contains every coset as a literal matrix set and all
its relevant intersections. The independent reader reconstructs the
entire displayed test and compares these sets, not just their counts.
Six damaged certificates are rejected; optimization modes are prohibited
because the historical acquisition code used assertions.

Every two-type flag is a translate of a base flag. It follows that the
incidence geometry is thin and that the stabilizer acts transitively
on each full panel. The argument of Lemma~\ref{actual} then gives flag
transitivity; its order-two panel actions are doubly transitive. Rank
three is a geometry because every edge has a common representative and
can be completed using that representative for the third type. The
residues of rank two are connected: their two order-two stabilizers
generate the displayed dihedral parabolic. The whole geometry is
connected because $G_0$ is maximal and $G_1\nleq G_0$, so
$\gen{G_0,G_1}=G$. This one marking consequently satisfies both readings
of local double transitivity. It yields no source correction and proves
no assertion about all markings.

\section{Attribution and scope}
The target is D.~Leemans, \emph{The rank 3 geometries of the simple
Suzuki groups $\Sz(q)$}, Note di Matematica 19(1) (1999), 43--63:
Theorem~5.3, its prime-field rank bound, and the rank-two and algebraic
local-action parts of the prime branches of Theorem~6.1. The structure
of Suzuki groups, their faithful matrices and ovoid action remain
classical inputs. The accepted L115 preliminary and normalizer-family
notes supply the affine complement argument, Frattini identity and
complete rank-two families with their equivalence counts.
Our derivation combines these through strict intersections and prime
subgroup covers; the matrix entry in Lemma~\ref{triangle} replaces one
implicit geometric step. No novelty claim is made for the elementary
lemmas. The one bounded computation supplies actual panels for an already
retained marking. No new conjecture credit, whole-source credit, Lean
verification or compression measurement is asserted.
\end{document}
